\documentclass[12pt]{article}
\usepackage[a4paper,margin=1.5cm]{geometry}
\usepackage[T1]{fontenc}
\usepackage[utf8]{inputenc}
\usepackage{lmodern}
\usepackage{amsmath,amssymb,booktabs,array}
\usepackage{longtable}
\usepackage{xcolor}
\usepackage{hyperref}

\hypersetup{
  colorlinks=true,
  linkcolor=blue,
  urlcolor=blue
}

\newcommand{\sgoverlay}[2]{\mathop{\mathchoice
  {\ooalign{\hfil$\displaystyle\sum$\hfil\cr\hfil$#1$\hfil\cr}}
  {\ooalign{\hfil$\textstyle\sum$\hfil\cr\hfil$#2$\hfil\cr}}
  {\ooalign{\hfil$\scriptstyle\sum$\hfil\cr\hfil$#2$\hfil\cr}}
  {\ooalign{\hfil$\scriptscriptstyle\sum$\hfil\cr\hfil$#2$\hfil\cr}}
}}

\newcommand{\sgsumintright}[1]{\mathop{\mathchoice
  {\ooalign{
    \hfil$\displaystyle\sum$\hfil\cr
    \hfil$\displaystyle\int$\hfil\cr
    \hidewidth\raisebox{-1.35ex}{\kern1.25ex$\scriptstyle #1$}\hidewidth\cr
  }}
  {\ooalign{
    \hfil$\textstyle\sum$\hfil\cr
    \hfil$\textstyle\int$\hfil\cr
    \hidewidth\raisebox{-1.00ex}{\kern1.00ex$\scriptscriptstyle #1$}\hidewidth\cr
  }}
  {\ooalign{
    \hfil$\scriptstyle\sum$\hfil\cr
    \hfil$\scriptstyle\int$\hfil\cr
    \hidewidth\raisebox{-0.75ex}{\kern0.75ex$\scriptscriptstyle #1$}\hidewidth\cr
  }}
  {\ooalign{
    \hfil$\scriptscriptstyle\sum$\hfil\cr
    \hfil$\scriptscriptstyle\int$\hfil\cr
    \hidewidth\raisebox{-0.60ex}{\kern0.55ex$\scriptscriptstyle #1$}\hidewidth\cr
  }}
}}

\newcommand{\sgdiamond}{\sgoverlay{\displaystyle\diamond}{\textstyle\diamond}}
\newcommand{\sgint}{\sgoverlay{\displaystyle\int}{\textstyle\int}}
\newcommand{\sgstar}{\sgoverlay{\displaystyle\star}{\textstyle\star}}
\newcommand{\sgcirc}{\sgoverlay{\displaystyle\circ}{\textstyle\circ}}
\newcommand{\sgsquare}{\sgoverlay{\displaystyle\square}{\textstyle\square}}
\newcommand{\sgtriangle}{\sgoverlay{\displaystyle\triangle}{\textstyle\triangle}}
\newcommand{\sgbowtie}{\sgoverlay{\displaystyle\bowtie}{\textstyle\bowtie}}
\newcommand{\sgbullet}{\sgoverlay{\displaystyle\bullet}{\textstyle\bullet}}
\newcommand{\sgtimes}{\sgoverlay{\displaystyle\times}{\textstyle\times}}
\newcommand{\sgplus}{\sgoverlay{\displaystyle+}{\textstyle+}}
\newcommand{\sgfinaldiamond}{\sgsumintright{\diamond}}
\newcommand{\sgfinalbullet}{\sgsumintright{\bullet}}
\newcommand{\sgfinalbowtie}{\sgsumintright{\bowtie}}

\newcommand{\sampleexpr}[1]{\[
  #1(X,M)=\sum_{i=1}^{N}\omega_i\bigl(a\mu_i+b\sigma_i+dE_i\bigr)
\]}

\begin{document}

\begin{center}
{\LARGE Sumigron Symbol Variants}\\[0.4em]
{\large Ten operator glyphs for visual selection}
\end{center}

\vspace{0.6em}

\noindent
This sheet is for choosing a visual operator sign for the Sumigron family. The strict notation can still remain
\[
\mathcal{S}_{\mathrm{sg}}(X,M)
\]
inside formal text, while a custom glyph may be used in diagrams, slides, or exploratory notes.

\section*{Recommended Reading Rule}
\begin{itemize}
\item If you want the sign to reflect your latest idea exactly, prefer \(\sgfinaldiamond\): sum plus integral with a tiny lower-right geometric marker.
\item If Sumigron is mainly about \textbf{structured aggregation}, prefer \(\sgdiamond\) or \(\sgbowtie\).
\item If it is mainly about \textbf{continuous accumulation}, prefer \(\sgint\).
\item If it is mainly about \textbf{salience and attention}, prefer \(\sgstar\) or \(\sgbullet\).
\item If it is mainly about \textbf{matrix/window blocks}, prefer \(\sgsquare\).
\end{itemize}

\section*{Preferred Hybrid}
Your current idea in strict form:
\[
\sgfinaldiamond(X,M)
\]

\noindent
Meaning:
\begin{itemize}
\item \(\sum\) gives aggregation,
\item \(\int\) gives continuous flow through the window,
\item the tiny lower-right marker gives Sumigron its own identity.
\end{itemize}

\noindent
Three close variants:
\[
\sgfinaldiamond(X,M), \qquad
\sgfinalbullet(X,M), \qquad
\sgfinalbowtie(X,M)
\]

\noindent
Interpretation:
\begin{itemize}
\item \(\sgfinaldiamond\): structured/geometric Sumigron
\item \(\sgfinalbullet\): compact attentive Sumigron
\item \(\sgfinalbowtie\): fusion-style Sumigron
\end{itemize}

\section*{Visual Catalog}
\renewcommand{\arraystretch}{1.3}
\begin{longtable}{>{\raggedright\arraybackslash}p{1.1cm} >{\centering\arraybackslash}p{2.0cm} >{\centering\arraybackslash}p{2.8cm} >{\raggedright\arraybackslash}p{8.3cm}}
\toprule
ID & Macro & Glyph & Intended reading \\
\midrule
\endhead
V1 & \texttt{\textbackslash sgfinaldiamond} & \(\sgfinaldiamond(X,M)\) & Sum plus integral with a tiny lower-right diamond. This is the direct implementation of your latest idea and the strongest candidate for a custom canonical Sumigron glyph. \\
V2 & \texttt{\textbackslash sgfinalbullet} & \(\sgfinalbullet(X,M)\) & Same hybrid base, but with a lower-right bullet. Cleaner and more compact if the diamond feels visually overloaded. \\
V3 & \texttt{\textbackslash sgfinalbowtie} & \(\sgfinalbowtie(X,M)\) & Same hybrid base, but with a lower-right bowtie. Better if you want the marker to imply fusion. \\
V4 & \texttt{\textbackslash sgdiamond} & \(\sgdiamond(X,M)\) & Sum plus diamond. Best if Sumigron is treated as a structured or geometric aggregator. This is the most semantically stable non-hybrid option. \\
V5 & \texttt{\textbackslash sgint} & \(\sgint(X,M)\) & Sum plus integral. Best if the operator should feel like a cumulative flow through the window rather than a plain reduction. \\
V6 & \texttt{\textbackslash sgstar} & \(\sgstar(X,M)\) & Sum plus star. Good for an attentive or salient variant, where peaks or informative fragments get emphasized. \\
V7 & \texttt{\textbackslash sgsquare} & \(\sgsquare(X,M)\) & Sum plus square. Good for matrix or block aggregation. Visually fits windows, tensors, and typed channel tables. \\
V8 & \texttt{\textbackslash sgbowtie} & \(\sgbowtie(X,M)\) & Sum plus bowtie. Best if Sumigron is interpreted as a fusion operator that mixes several channels into one structured score. \\
V9 & \texttt{\textbackslash sgbullet} & \(\sgbullet(X,M)\) & Sum plus bullet. Minimal attentive marker. Strong if the operator should feel compact and focused. \\
V10 & \texttt{\textbackslash sgtimes} & \(\sgtimes(X,M)\) & Sum plus cross. Reads like interaction, coupling, or cross-channel synthesis. Useful if nonlinear interaction is central. \\
\bottomrule
\end{longtable}

\section*{In Formula Form}
Below the same operators are rendered in the same Sumigron expression template.

\sampleexpr{\sgfinaldiamond}
\sampleexpr{\sgfinalbullet}
\sampleexpr{\sgfinalbowtie}
\sampleexpr{\sgdiamond}
\sampleexpr{\sgint}
\sampleexpr{\sgstar}
\sampleexpr{\sgsquare}
\sampleexpr{\sgbowtie}
\sampleexpr{\sgbullet}
\sampleexpr{\sgtimes}

\section*{Pragmatic Choice}
If one sign should be chosen right now:
\begin{itemize}
\item \textbf{Best custom canonical sign}: \(\sgfinaldiamond\)
\item \textbf{Best fallback if the diamond is too busy}: \(\sgfinalbullet\)
\item \textbf{Best formal default without custom corner marker}: \(\sgdiamond\)
\item \textbf{Best for typed matrix aggregation}: \(\sgsquare\)
\item \textbf{Best for multimodal fusion}: \(\sgbowtie\)
\end{itemize}

\noindent
The clean fallback notation remains
\[
\mathcal{S}_{\mathrm{sg}}(X,M), \qquad
\mathcal{S}^{\mathrm{att}}_{\mathrm{sg}}(X,M).
\]

\end{document}
